\documentclass[11pt,letterpaper]{article}
\usepackage[utf8]{inputenc}
\usepackage[spanish,es-tabla]{babel}
\usepackage[margin=2.5cm]{geometry}
\usepackage{amsmath,amssymb}
\usepackage{booktabs}
\usepackage{graphicx}
\usepackage{float}
\usepackage{fancyhdr}
\usepackage{hyperref}
\usepackage[table]{xcolor}
\usepackage{tikz}
\usepackage{transparent}
\usepackage{eso-pic}
\usepackage{caption}
\definecolor{GeoOrange}{HTML}{FF8C00}
\definecolor{GeoBlue}{HTML}{0D47A1}
\AddToShipoutPictureBG{
  \AtPageCenter{
    \makebox[0pt]{\transparent{0.06}
\includegraphics[width=12cm]{Iconos/logo_geopol.png}}
  }
}
\hypersetup{colorlinks=true, linkcolor=GeoBlue, urlcolor=GeoOrange}
\pagestyle{fancy}
\fancyhf{}
\fancyhead[L]{\footnotesize \textcolor{GeoBlue}{\textbf{GeoPol Web}} - Reporte Técnico}
\fancyhead[R]{\footnotesize Universidad Distrital F.J.C.}
\fancyfoot[C]{\thepage}
\renewcommand{\headrulewidth}{0.4pt}
\renewcommand{\footrulewidth}{0.4pt}
\begin{document}
\begin{titlepage}
\begin{tikzpicture}[remember picture,overlay]
  \fill[GeoBlue] (current page.north west) rectangle ([yshift=-4cm]current page.north east);
  \fill[GeoOrange] ([yshift=-4cm]current page.north west) rectangle ([yshift=-4.5cm]current page.north east);
  \fill[GeoBlue!5] (current page.south west) rectangle ([yshift=4cm]current page.south east);
\end{tikzpicture}
\vspace*{-2cm}
\begin{center}
  \textcolor{white}{\Huge \textbf{GEOPORTAL WEB (GeoPol)}} \\[0.5cm]
  \textcolor{white}{\large \textit{''Máxima precisión al alcance de tus manos''}} \\[2cm]
  \vspace{3cm}
  {\LARGE \textbf{INFORME TÉCNICO DE INGENIERÍA}} \\[0.5cm]
  {\Large \textbf{Abierta}} \\[2cm]
  \begin{flushleft}
    \Large \textbf{Autores del Procesamiento:}\\[0.2cm]
    \large $\bullet$ Kevin Stiven Cubillos Ramirez \\[0.1cm]
    \large $\bullet$ Sergio Eduardo Barbosa Torres \\[0.1cm]
    \vspace{1cm}
    \Large \textbf{Tutor - Director de Proyecto:}\\[0.2cm]
    \large $\bullet$ Ing. Edgar Ladino
  \end{flushleft}
  \vfill
  \textbf{UNIVERSIDAD DISTRITAL FRANCISCO JOSÉ DE CALDAS}\\[0.2cm]
  Facultad de Medio Ambiente y Recursos Naturales \\ Ingeniería Topográfica / Civil \\[0.2cm]
  Bogotá D.C. -- 28 de July de 2026
\end{center}
\end{titlepage}
\tableofcontents
\newpage
\section{Marco Teórico y Referencia Geodésica}
El presente informe detalla el cálculo, ajuste y representación de una red de apoyo planimétrico. Este procesamiento se apoya en los lineamientos técnicos de la Topografía Clásica y la normatividad geodésica dictada por el \textbf{Instituto Geográfico Agustín Codazzi (IGAC)}.
\subsection{Sistema de Georreferenciación (MAGNA-SIRGAS)}
De acuerdo con la \textbf{Resolución 471 de 2020} (y su modificación en la Resolución 529 de 2020) emitida por el IGAC, el único sistema oficial de coordenadas para la República de Colombia es el Origen Nacional \textbf{MAGNA-SIRGAS (EPSG: 9377)}.
\subsection{Fundamento Matemático: Abierta}
Una Poligonal Abierta con Control es aquella que parte de una línea base con azimut y coordenadas conocidas, y finaliza su recorrido en otra estación (o par de estaciones) de coordenadas igualmente conocidas.
\begin{itemize}
  \item \textbf{Cierre Angular:} El azimut calculado del último alineamiento se compara contra el azimut teórico conocido de llegada.
  \item \textbf{Cierre Lineal:} La sumatoria de las proyecciones debe ser igual a la diferencia real de coordenadas entre el punto de llegada y el de partida ($\Sigma \Delta N = N_{llegada} - N_{partida}$). Los desvíos se compensan proporcionalmente a las distancias.
\end{itemize}
\section{Trabajo de Campo: Registro de Observaciones}
En la siguiente sección se relacionan los datos brutos levantados en campo (ángulos horizontales, verticales, distancias inclinadas y alturas instrumentales).
\begin{table}[H]
  \centering
  \small
  \caption{Cartera de Observaciones Brutas (Parte 1)}
  \label{tab:campo_0}
  \begin{tabular}{|c|c|c|c|c|c|}
    \hline
    \rowcolor{GeoBlue}
    \textcolor{white}{\textbf{Estacionado}} & \textcolor{white}{\textbf{Pto\_Obs}} & \textcolor{white}{\textbf{Hz\_G}} & \textcolor{white}{\textbf{Hz\_M}} & \textcolor{white}{\textbf{Hz\_S}} & \textcolor{white}{\textbf{Z\_G}} \\
    \hline
    \rowcolor{GeoBlue!5}
    GPS-09 & C-1 & 76.000 & 56.000 & 32.000 & 81.000 \\
    \hline
    \rowcolor{white}
    C-1 & C-2 & 246.000 & 58.000 & 41.000 & 89.000 \\
    \hline
    \rowcolor{GeoBlue!5}
    C-2 & C-3 & 135.000 & 51.000 & 53.000 & 89.000 \\
    \hline
    \rowcolor{white}
    C-3 & C-4 & 223.000 & 25.000 & 11.000 & 89.000 \\
    \hline
    \rowcolor{GeoBlue!5}
    C-4 & C-5 & 205.000 & 20.000 & 14.000 & 90.000 \\
    \hline
    \rowcolor{white}
    C-5 & C-6 & 197.000 & 17.000 & 36.000 & 89.000 \\
    \hline
    \rowcolor{GeoBlue!5}
    C-6 & C-7 & 162.000 & 49.000 & 57.000 & 90.000 \\
    \hline
    \rowcolor{white}
    C-7 & GPS-06 & 180.000 & 0.000 & 0.000 & 90.000 \\
    \hline
    \rowcolor{GeoBlue!5}
    GPS-06 & GPS-05 & 180.000 & 0.000 & 0.000 & 90.000 \\
    \hline
  \end{tabular}
\end{table}
\vspace{0.8cm}
\begin{table}[H]
  \centering
  \small
  \caption{Cartera de Observaciones Brutas (Parte 2)}
  \label{tab:campo_1}
  \begin{tabular}{|c|c|c|c|c|c|}
    \hline
    \rowcolor{GeoBlue}
    \textcolor{white}{\textbf{Estacionado}} & \textcolor{white}{\textbf{Pto\_Obs}} & \textcolor{white}{\textbf{Z\_M}} & \textcolor{white}{\textbf{Z\_S}} & \textcolor{white}{\textbf{Dist\_Inc}} & \textcolor{white}{\textbf{hi}} \\
    \hline
    \rowcolor{GeoBlue!5}
    GPS-09 & C-1 & 4.000 & 20.000 & 20.119 & 1.398 \\
    \hline
    \rowcolor{white}
    C-1 & C-2 & 45.000 & 10.000 & 73.699 & 1.470 \\
    \hline
    \rowcolor{GeoBlue!5}
    C-2 & C-3 & 30.000 & 13.000 & 116.226 & 1.528 \\
    \hline
    \rowcolor{white}
    C-3 & C-4 & 39.000 & 21.000 & 96.228 & 1.537 \\
    \hline
    \rowcolor{GeoBlue!5}
    C-4 & C-5 & 14.000 & 0.000 & 47.085 & 1.534 \\
    \hline
    \rowcolor{white}
    C-5 & C-6 & 43.000 & 33.000 & 32.462 & 1.563 \\
    \hline
    \rowcolor{GeoBlue!5}
    C-6 & C-7 & 33.000 & 19.000 & 58.209 & 1.550 \\
    \hline
    \rowcolor{white}
    C-7 & GPS-06 & 0.000 & 0.000 & 50.000 & 1.500 \\
    \hline
    \rowcolor{GeoBlue!5}
    GPS-06 & GPS-05 & 0.000 & 0.000 & 50.000 & 1.500 \\
    \hline
  \end{tabular}
\end{table}
\vspace{0.8cm}
\begin{table}[H]
  \centering
  \small
  \caption{Cartera de Observaciones Brutas (Parte 3)}
  \label{tab:campo_2}
  \begin{tabular}{|c|c|c|}
    \hline
    \rowcolor{GeoBlue}
    \textcolor{white}{\textbf{Estacionado}} & \textcolor{white}{\textbf{Pto\_Obs}} & \textcolor{white}{\textbf{hr}} \\
    \hline
    \rowcolor{GeoBlue!5}
    GPS-09 & C-1 & 1.800 \\
    \hline
    \rowcolor{white}
    C-1 & C-2 & 1.800 \\
    \hline
    \rowcolor{GeoBlue!5}
    C-2 & C-3 & 1.800 \\
    \hline
    \rowcolor{white}
    C-3 & C-4 & 1.800 \\
    \hline
    \rowcolor{GeoBlue!5}
    C-4 & C-5 & 1.800 \\
    \hline
    \rowcolor{white}
    C-5 & C-6 & 1.800 \\
    \hline
    \rowcolor{GeoBlue!5}
    C-6 & C-7 & 1.800 \\
    \hline
    \rowcolor{white}
    C-7 & GPS-06 & 1.800 \\
    \hline
    \rowcolor{GeoBlue!5}
    GPS-06 & GPS-05 & 1.800 \\
    \hline
  \end{tabular}
\end{table}
\vspace{0.8cm}

\section{Cálculo, Análisis de Errores y Compensación}
El motor de cálculo topográfico procesó las observaciones, obteniendo las siguientes métricas de cierre (incluyendo planimetría y altimetría) antes de ejecutar el ajuste perimetral:
\begin{table}[H]
  \centering
  \caption{Métricas de Cierre Geométrico}
  \begin{tabular}{|l|l|}
    \hline
    \rowcolor{GeoBlue}
    \textcolor{white}{\textbf{Parámetro Analizado}} & \textcolor{white}{\textbf{Magnitud del Error}} \\
    \hline
    Error Angular Bruto & -137.4027634615929 \\
    \rowcolor{GeoBlue!5}
    Error Horizontal Este ($e_x$) & -86.19581 m \\
    Error Horizontal Norte ($e_y$) & 268.39489 m \\
    \rowcolor{white}
    Error Vertical ($\Delta Z$) & 2.27667 m \\
    Precisión Vertical Relativa & 1 en 216 \\
    \rowcolor{GeoBlue!5}
    Error Lineal Cierre ($e_L$) & 281.89632 m \\
    Precisión Planimétrica ($1:P$) & 1 en 1 \\
    \hline
  \end{tabular}
\end{table}
\subsection{Cartera Final de Coordenadas Ajustadas}
\begin{table}[H]
  \centering
  \small
  \caption{Coordenadas Compensadas de la Red (Parte 1)}
  \label{tab:ajuste_0}
  \begin{tabular}{|c|c|c|c|c|c|}
    \hline
    \rowcolor{GeoBlue}
    \textcolor{white}{\textbf{Estacionado}} & \textcolor{white}{\textbf{Pto\_Obs}} & \textcolor{white}{\textbf{Ang\_Hz\_Ajus}} & \textcolor{white}{\textbf{Azimut\_Línea}} & \textcolor{white}{\textbf{Dep\_Ajus (E)}} & \textcolor{white}{\textbf{Lat\_Ajus (N)}} \\
    \hline
    \rowcolor{GeoBlue!5}
    GPS-09 & C-1 & 92° 12' 33.11" & 173° 4' 59.39" & 5.863 & -30.534 \\
    \hline
    \rowcolor{white}
    C-1 & C-2 & 262° 14' 42.11" & 255° 19' 41.50" & -58.430 & -58.726 \\
    \hline
    \rowcolor{GeoBlue!5}
    C-2 & C-3 & 151° 7' 54.11" & 226° 27' 35.60" & -63.960 & -143.234 \\
    \hline
    \rowcolor{white}
    C-3 & C-4 & 238° 41' 12.11" & 285° 8' 47.71" & -76.086 & -27.162 \\
    \hline
    \rowcolor{GeoBlue!5}
    C-4 & C-5 & 220° 36' 15.11" & 325° 45' 2.81" & -18.280 & 13.327 \\
    \hline
    \rowcolor{white}
    C-5 & C-6 & 212° 33' 37.11" & 358° 18' 39.92" & 4.710 & 14.803 \\
    \hline
    \rowcolor{GeoBlue!5}
    C-6 & C-7 & 178° 5' 58.11" & 356° 24' 38.02" & 6.517 & 26.454 \\
    \hline
    \rowcolor{white}
    C-7 & GPS-06 & 195° 16' 1.11" & 11° 40' 39.13" & 18.848 & 21.787 \\
    \hline
    \rowcolor{GeoBlue!5}
    GPS-06 & GPS-05 & 195° 16' 1.11" & 26° 56' 40.23" & 0.000 & 0.000 \\
    \hline
  \end{tabular}
\end{table}
\vspace{0.8cm}
\begin{table}[H]
  \centering
  \small
  \caption{Coordenadas Compensadas de la Red (Parte 2)}
  \label{tab:ajuste_1}
  \begin{tabular}{|c|c|c|c|c|}
    \hline
    \rowcolor{GeoBlue}
    \textcolor{white}{\textbf{Estacionado}} & \textcolor{white}{\textbf{Pto\_Obs}} & \textcolor{white}{\textbf{X\_Estacion}} & \textcolor{white}{\textbf{Y\_Estacion}} & \textcolor{white}{\textbf{Z\_Estacion}} \\
    \hline
    \rowcolor{GeoBlue!5}
    GPS-09 & C-1 & 86138.390 & 102562.748 & 2565.979 \\
    \hline
    \rowcolor{white}
    C-1 & C-2 & 86144.253 & 102532.214 & 2568.608 \\
    \hline
    \rowcolor{GeoBlue!5}
    C-2 & C-3 & 86085.823 & 102473.488 & 2568.256 \\
    \hline
    \rowcolor{white}
    C-3 & C-4 & 86021.863 & 102330.254 & 2568.455 \\
    \hline
    \rowcolor{GeoBlue!5}
    C-4 & C-5 & 85945.778 & 102303.093 & 2568.326 \\
    \hline
    \rowcolor{white}
    C-5 & C-6 & 85927.498 & 102316.419 & 2567.651 \\
    \hline
    \rowcolor{GeoBlue!5}
    C-6 & C-7 & 85932.208 & 102331.222 & 2567.420 \\
    \hline
    \rowcolor{white}
    C-7 & GPS-06 & 85938.725 & 102357.676 & 2566.338 \\
    \hline
    \rowcolor{GeoBlue!5}
    GPS-06 & GPS-05 & 85957.573 & 102379.463 & 2565.807 \\
    \hline
  \end{tabular}
\end{table}
\vspace{0.8cm}

\section{Esquema Geométrico de la Red Planimétrica}
\begin{figure}[H]
  \centering
  \includegraphics[width=0.95\textwidth]{Plano_Exportado.png}
  \caption{Plano As-Built de la Abierta}
\end{figure}
\section{Conclusiones y Dictamen Técnico}
\begin{itemize}
  \item Se obtuvo una precisión de 1:1. Esta precisión es \textbf{DEFICIENTE}. No cumple con los estándares mínimos para levantamientos topográficos convencionales. Se recomienda enfáticamente revisar los ángulos observados o repetir la medición en campo.
  \item El procesamiento fue realizado exitosamente de forma automatizada mediante \textbf{GeoPol Web}, suprimiendo el error de cálculo humano y garantizando la trazabilidad matemática requerida por la ingeniería civil moderna.
\end{itemize}
\end{document}